% eksperymentalne tłumaczenie na włoski

%

\subsection{Ortocentrum i trzy wysokości}

\begin{definition}[wysokość]
\index{wysokość trójkąta}%
    \pol{Niech $\triangle ABC$ będzie trójkątem, w którym przez punkt $C$ poprowadzono prostą prostopadłą do boku $AB$.}
    \ita{Sia $\triangle ABC$ un triangolo e si tracci per il punto $C$ la retta perpendicolare al lato $AB$.}
    \pol{Odcinek leżący na tej prostej, którego jeden koniec znajduje się w punkcie $C$, zaś drugi leży na boku $AB$, nazywamy wysokością opuszczoną z punktu $C$, drugi jego koniec to spodek wysokości.}
    \ita{Il segmento su tale retta avente un estremo nel punto $C$ e l'altro sul lato $AB$ si chiama altezza relativa al punto $C$, e il secondo estremo è il piede dell'altezza.}
\end{definition}

\pol{Po angielsku mamy \emph{altitude} oraz \emph{foot of the altitude}.}
\ita{In inglese si usano i termini \emph{altitude} e \emph{foot of the altitude}.}

% TODO: Sometimes an arbitrary edge is chosen to be the base, in which case the opposite vertex is called the apex; the shortest segment between the base and apex is the height. The area of a triangle equals one-half the product of height and base length.

\begin{proposition}
\label{wysokosci_przecinaja_sie}%
	\pol{Trzy wysokości trójkąta przecinają się w jednym punkcie zwanym ortocentrum.}
	\ita{Le tre altezze di un triangolo si intersecano in un unico punto chiamato ortocentro.}
\index{ortocentrum}%
\end{proposition}

\pol{Wysokości trójkąta rozwartokątnego nie przecinają się, musimy je więc zastąpić przez proste, na których leżą.}
\ita{Le altezze di un triangolo ottusangolo non si intersecano all'interno del triangolo, quindi si considerano le rette su cui giacciono.}

\pol{Dowodów we współczesnej literaturze nie brakuje: są u Pompego \cite[s. 38]{pompe_2022}, gdzie zmyślnie używa równoległoboków, Guzickiego \cite[s. 218]{guzicki_2021}, Zetela \cite[s. 14, 151]{zetel_2020}.}
\ita{Le dimostrazioni nella letteratura moderna non mancano: si trovano in Pompe \cite[p.~38]{pompe_2022}, dove si usano in modo ingegnoso i parallelogrammi, in Guzicki \cite[p.~218]{guzicki_2021}, e in Zetel \cite[pp.~14, 151]{zetel_2020}.}
\pol{Warto też przeczytać tekst Hartshorne'a \cite[s. 54]{hartshorne2000} oraz Audina \cite[s. 61]{audin_2003}.}
\ita{Vale anche la pena leggere Hartshorne \cite[p.~54]{hartshorne2000} e Audin \cite[p.~61]{audin_2003}.}

\pol{Ortocentrum (\emph{orthocenter}) będzie czasami nazywane środkiem ortycznym, ale nazwa ta nie utrzyma się długo.}
\ita{L'ortocentro (\emph{orthocenter}) viene talvolta chiamato centro ortico, ma questo termine non si affermerà a lungo.}

\begin{proposition}
    \pol{Dany jest trójkąt $ABC$ oraz jego ortocentrum $H$.}
    \ita{Sia dato un triangolo $ABC$ e il suo ortocentro $H$.}
    \pol{Odbicia symetryczne punktu $H$ względem prostych zawierających boki trójkąta $ABC$ leżą na okręgu opisanym na tym trójkącie.}
    \ita{Le simmetrie del punto $H$ rispetto alle rette contenenti i lati del triangolo $ABC$ giacciono sulla circonferenza circoscritta al triangolo.}
\end{proposition}

\pol{Dowód tego prostego stwierdzenia przedstawia Neugebauer \cite[s. 28]{neugebauer_2018}.}
\ita{La dimostrazione di questo semplice fatto è presentata in Neugebauer \cite[p.~28]{neugebauer_2018}.}

%